\scanpage{8}
\mspage{6}
\inserted{Mais} Soit $k$ un anneau comm (ass. unit.) et $T$ une \unsure{indéterminé|indéterminée}, on prendra dorénavant
\[
\begin{aligned}
S&=k\{T\},\qquad J=k\{T\}^{+}=\operatorname{Ker}(k\{T\}\unsure{\text{---}|\to} k)\\
t&=T
\end{aligned}
\]
Donc
\[
\left\{\begin{aligned}
C^\bullet_{\mathrm{DRpd}\,[n]}(S,J,t)&\simeq
\underbrace{k\{T,X_0,\ldots,X_n\}/(\textstyle\sum X_i-T)_{\mathrm{pd}}}_{\substack{\wr\\ k\{X_0,\ldots,X_n\}}}
\otimes_k\bigwedge^{\bullet}k^{[n]}/k\\
S/J&\simeq k
\end{aligned}\right.
\]
Soit
\[
\begin{aligned}
\Phi_{k\ast}&=\big[[n]\longmapsto k^{[n]}\big]\quad\underset{\substack{\text{inclusion}\\\text{diagonale}}}{\supset}\ k_\ast=\big[[n]\text{---}k\big]\\
\Psi_{k\ast}&=\Phi_{k\ast}/\unsure{(k_\ast)|k_\ast}=\big([n]\text{---}k^{[n]}/\underbrace{k}_{\mathrm{diag}}\big)
\end{aligned}
\]
On a alors
\[
\begin{aligned}
C^\bullet_{\mathrm{DRpd}\,\ast\deleted{\unsure{n}}}(k\{T\},k\{T\}^{+},k)&\simeq\unsure{\Gamma|\mathscr T}^{\bullet}_{k}\Phi_{k\ast}\otimes_k\bigwedge^\bullet\Psi_{k\ast}\\
\underset{\text{déf}}{\Vert}&\\[-2pt]
C^{\bullet\bullet}_{\mathrm{DRpd}\,\ast}(k)&
\end{aligned}
\]
\[
\boxed{C^{\bullet\bullet}_{\mathrm{DRpd}\,\ast}(k)\simeq\Gamma^\bullet_k\Phi_{k\ast}\otimes_k\bigwedge^\bullet\Psi_{k\ast}}
\]

\uline{Structure}. \inserted{$k\{T\}$-}Algèbre \deleted{\unsure{diff.}} \uline{bigraduée} (deg. \unsure{complém.} \uline{et} degré extérieur, d'\uline{où} degré total) \inserted{unit.} assoc. \deleted{anticomm.}\inserted{alternée} (anticomm + carré des éléments de degré impair nuls), augmentée vers $k$, à p.d. sur l'idéal d'augm, avec diff de bidegré $-1,+1$ (donc degré total 0) \inserted{compatible avec p.d.}
% slanted
\margin{la bigraduation provient de la \uline{graduation} de $S$, et du fait que $t$ est homogène (ici de degré 1)}

% archive sheet number: 7
